% ============================================================
% Les fantômes d'Ombrequatre — Recapitulatif solveur C + base de donnees
% Copier-coller ce fichier sur Overleaf (compiler avec pdfLaTeX)
% ============================================================
\documentclass[11pt,a4paper]{article}

\usepackage[T1]{fontenc}
\usepackage[utf8]{inputenc}
\usepackage[french]{babel}
\usepackage{geometry}
\geometry{margin=2.2cm}
\usepackage{graphicx}
\usepackage{booktabs}
\usepackage{longtable}
\usepackage{array}
\usepackage{hyperref}
\usepackage{xcolor}
\usepackage{listings}
\usepackage{tikz}
\usetikzlibrary{arrows.meta, positioning, shapes.geometric}

\hypersetup{
    colorlinks=true,
    linkcolor=blue!60!black,
    urlcolor=blue!60!black,
    citecolor=blue!60!black
}

\definecolor{codebg}{RGB}{245,245,250}
\definecolor{keyword}{RGB}{0,80,160}

\lstdefinestyle{kllisting}{
    backgroundcolor=\color{codebg},
    basicstyle=\ttfamily\small,
    breaklines=true,
    frame=single,
    rulecolor=\color{gray!40},
    columns=fullflexible,
    keepspaces=true,
    showstringspaces=false,
    literate={é}{{\'e}}1 {è}{{\`e}}1 {ê}{{\^e}}1 {à}{{\`a}}1
             {ù}{{\`u}}1 {ô}{{\^o}}1 {î}{{\^i}}1 {ç}{{\c{c}}}1
             {É}{{\'E}}1 {ë}{{\"e}}1
}

\lstdefinestyle{kllistingNoFrame}{
    backgroundcolor=\color{codebg},
    basicstyle=\ttfamily\small,
    breaklines=true,
    columns=fullflexible,
    keepspaces=true,
    showstringspaces=false
}

\title{%
  \textbf{Les fantômes d'Ombrequatre}\\[0.4em]
  \large Recapitulatif detaille --- Solveur C, base de donnees\\
  \large et export de solutions en fichiers \texttt{.txt}
}
\author{CIR1 2025--2026 --- Groupe 4doigtsdelamain}
\date{Mai 2026}

\begin{document}
\maketitle
\tableofcontents
\newpage

%------------------------------------------------------------
\section{Objectif du travail realise}
%------------------------------------------------------------

Ce document recapitule l'ensemble des developpements realises pour repondre a la demande suivante :

\begin{quote}
\textit{Proposer et coder une solution pour le solveur en C prenant en compte les niveaux a partir de la base de donnees, et renvoyer la solution sous la forme d'un fichier \texttt{.txt} pour traitement ulterieur.}
\end{quote}

\textbf{Points cles :}
\begin{itemize}
    \item Les niveaux sont lus depuis MySQL (table \texttt{niveau}, colonne \texttt{map}).
    \item Le calcul de la solution est effectue \textbf{exclusivement en C} (BFS existant dans \texttt{solver/solver.c}).
    \item Chaque niveau produit un fichier texte structure dans \texttt{output/solutions/}.
    \item Le solveur JavaScript (\texttt{solver-worker.js}) et le site web \textbf{n'ont pas ete modifies} dans cette etape : il s'agit d'un pipeline \textbf{parallele}, pret pour import futur en BDD.
\end{itemize}

%------------------------------------------------------------
\section{Contexte du projet}
%------------------------------------------------------------

\subsection{Jeu et mecanique modelisee}

Les fantômes d'Ombrequatre est un jeu de labyrinthe ou le chevalier :
\begin{itemize}
    \item \textbf{glisse} dans une direction (U/D/L/R) jusqu'a un mur ou un croisement ;
    \item compte \textbf{un coup} par glissement (quelle que soit la longueur) ;
    \item \textbf{ramasse} les gemmes sur le trajet (\texttt{.}, \texttt{o}, \texttt{c}) ;
    \item ne peut \textbf{pas} choisir la direction opposee au coup precedent.
\end{itemize}

\textbf{Objectif du solveur :} trouver la sequence de glissements de \textbf{longueur minimale} qui ramasse toutes les gemmes (BFS = parcours en largeur optimal).

\subsection{Etat du projet avant cette evolution}

\begin{table}[h]
\centering
\small
\begin{tabular}{@{}ll@{}}
\toprule
\textbf{Composant} & \textbf{Role} \\
\midrule
\texttt{solver/solver.c} & BFS C : \texttt{kl\_solve()} (gems-only), \texttt{kl\_solve\_safe()} (R/G/B) \\
\texttt{solver/main.c} & CLI unitaire (\texttt{./solver level.txt}) \\
\texttt{web/js/solver-worker.js} & BFS en navigateur (Worker), jaune inclus \\
\texttt{web/api/get\_solution.php} & Lit \texttt{solution\_cache} en BDD \\
Table \texttt{niveau} & \texttt{map}, \texttt{solution\_cache}, \texttt{solution\_safe} \\
\bottomrule
\end{tabular}
\caption{Composants existants}
\end{table}

%------------------------------------------------------------
\section{Architecture retenue}
%------------------------------------------------------------

Plutot que de lier le programme C directement a MySQL (\texttt{libmysqlclient}, compilation lourde sous Windows/MAMP), l'architecture adoptee est en \textbf{deux etapes} :

\begin{enumerate}
    \item \textbf{PHP CLI} --- connexion PDO, export des cartes et creation d'un manifeste.
    \item \textbf{C (\texttt{solve\_db})} --- lecture du manifeste, resolution, ecriture des fichiers solution.
\end{enumerate}

Cette approche correspond a la \textbf{voie A} du document d'architecture (PHP orchestre, C calcule).

\begin{figure}[h]
\centering
\begin{tikzpicture}[
    node distance=1.1cm and 1.8cm,
    box/.style={draw, rounded corners, minimum width=2.6cm, minimum height=0.9cm, align=center, font=\small},
    db/.style={box, fill=blue!8},
    php/.style={box, fill=green!10},
    file/.style={box, fill=orange!10},
    c/.style={box, fill=red!8}
]
    \node[db] (mysql) {MySQL\\\texttt{niveau}};
    \node[php, right=of mysql] (php) {\texttt{solveurCrenvoyanttxt/}\\\texttt{cli/solve\_from\_db.php}};
    \node[file, below=0.9cm of php, xshift=-1.4cm] (levels) {\texttt{.../output/levels/}\\\texttt{level\_N.txt}};
    \node[file, below=0.9cm of php, xshift=1.4cm] (manifest) {\texttt{levels\_manifest.txt}};
    \node[c, right=3.2cm of php] (solve) {\texttt{solve\_db.exe}};
    \node[file, right=of solve] (sol) {\texttt{.../output/solutions/}\\\texttt{level\_N\_solution.txt}};

    \draw[-{Stealth}] (mysql) -- node[above, font=\scriptsize] {SELECT id, map} (php);
    \draw[-{Stealth}] (php) -- (levels);
    \draw[-{Stealth}] (php) -- (manifest);
    \draw[-{Stealth}] (manifest) -- (solve);
    \draw[-{Stealth}] (levels) -- node[above, font=\scriptsize, pos=0.35] {chemins} (solve);
    \draw[-{Stealth}] (solve) -- (sol);
\end{tikzpicture}
\caption{Pipeline export BDD $\rightarrow$ resolution C $\rightarrow$ fichiers \texttt{.txt}}
\end{figure}

%------------------------------------------------------------
\section{Fichiers crees}
%------------------------------------------------------------

\subsection{\texttt{solveurCrenvoyanttxt/src/batch\_io.h} et \texttt{batch\_io.c}}

Couche utilitaire partagee entre le solveur unitaire (\texttt{solver}) et le batch (\texttt{solve\_db}).

\begin{table}[h]
\centering
\small
\begin{tabular}{@{}p{4.2cm}p{9.5cm}@{}}
\toprule
\textbf{Fonction} & \textbf{Description} \\
\midrule
\texttt{kl\_sequence\_string()} & Convertit \texttt{kl\_solution\_t} en chaine compacte \texttt{UDLR} \\
\texttt{kl\_solve\_level\_file()} & Parse, resout (safe-path + fallback optionnel) \\
\texttt{kl\_write\_solution\_txt()} & Ecrit le fichier solution (format cle=valeur) \\
\bottomrule
\end{tabular}
\end{table}

\textbf{Logique de \texttt{kl\_solve\_level\_file()} :}
\begin{enumerate}
    \item \texttt{kl\_parse\_level\_file()} lit le fichier au format projet.
    \item Si \texttt{use\_safe} : appel \texttt{kl\_solve\_safe()} (fantomes R/G/B ; \textbf{jaune ignore}).
    \item Si echec et \texttt{allow\_fallback} : appel \texttt{kl\_solve()} (gems-only).
    \item Renseigne \texttt{mode\_out} : \texttt{"safe-path"} ou \texttt{"gems-only"}.
\end{enumerate}

\subsection{\texttt{solveurCrenvoyanttxt/src/solve\_db.c}}

Nouveau programme batch. Binaire produit : \texttt{solve\_db} ou \texttt{solve\_db.exe}.

\textbf{Ligne de commande :}
\begin{lstlisting}[style=kllisting]
solve_db --manifest <fichier> --output-dir <dossier> [options]
\end{lstlisting}

\begin{table}[h]
\centering
\small
\begin{tabular}{@{}ll@{}}
\toprule
\textbf{Option} & \textbf{Effet} \\
\midrule
\texttt{--manifest} & Fichier : une ligne \texttt{id|chemin/level\_N.txt} par niveau \\
\texttt{--output-dir} & Dossier de sortie (\texttt{level\_<id>\_solution.txt}) \\
\texttt{--safe-path} & BFS evitant R/G/B (defaut) \\
\texttt{--gems-only} & BFS sans fantomes \\
\texttt{--fallback} & Si safe-path echoue, retenter gems-only \\
\bottomrule
\end{tabular}
\end{table}

\textbf{Comportement :} creation du dossier de sortie, lecture du manifeste, resolution niveau par niveau, resume console (OK / UNSOLVABLE / erreur), code retour 0 ou 1.

\subsection{\texttt{solveurCrenvoyanttxt/cli/solve\_from\_db.php}}

Point d'entree depuis la base de donnees.

\begin{enumerate}
    \item Charge \texttt{web/includes/config.php}.
    \item Connexion PDO directe (hors \texttt{getDB()} HTTP).
    \item \texttt{SELECT id, map FROM niveau ORDER BY id}.
    \item Ecrit \texttt{output/levels/level\_<id>.txt} pour chaque niveau.
    \item Ecrit \texttt{output/levels\_manifest.txt}.
    \item Lance \texttt{solve\_db} avec \texttt{exec()}.
\end{enumerate}

\textbf{Options PHP :}
\begin{itemize}
    \item \texttt{--gems-only} : passe \texttt{--gems-only} au binaire C.
    \item \texttt{--no-fallback} : pas de repli gems-only.
\end{itemize}

\subsection{Dossiers et documentation}

\begin{itemize}
    \item \texttt{solveurCrenvoyanttxt/output/levels/}, \texttt{output/solutions/}
    \item \texttt{solveurCrenvoyanttxt/build.bat}, \texttt{Makefile}
    \item \texttt{solveurCrenvoyanttxt/README.md} --- mode d'emploi
\end{itemize}

%------------------------------------------------------------
\section{Fichiers modifies}
%------------------------------------------------------------

\subsection{\texttt{solveurCrenvoyanttxt/src/solve\_one.c}}

CLI pour un seul niveau $\rightarrow$ fichier \texttt{.txt} :
\begin{lstlisting}[style=kllisting]
solve_one output/levels/level_1.txt --output output/solutions/level_1_solution.txt --safe-path --fallback
\end{lstlisting}

\subsection{\texttt{solveurCrenvoyanttxt/Makefile}}

Lie le moteur partage \texttt{../solver/solver.c}. Cibles : \texttt{solve\_db}, \texttt{solve\_one}.

Le dossier \texttt{solver/} du projet reste le CLI unitaire d'origine (\texttt{solver + main.c}), sans batch.

%------------------------------------------------------------
\section{Formats de donnees}
%------------------------------------------------------------

\subsection{Manifeste (\texttt{output/levels\_manifest.txt})}

Une ligne par niveau :
\begin{lstlisting}[style=kllistingNoFrame]
1|c:\MAMP\htdocs\project\output\levels\level_1.txt
2|c:\MAMP\htdocs\project\output\levels\level_2.txt
\end{lstlisting}

Separateur : \texttt{|}. Lignes vides et commentaires \texttt{\#} ignores.

\subsection{Fichier solution (\texttt{level\_<id>\_solution.txt})}

\begin{lstlisting}[style=kllisting]
# Les fantômes d'Ombrequatre - solution export
level_id=4
level_file=C:\MAMP\htdocs\project\output\levels\level_4.txt
solvable=yes
mode=safe-path
move_count=15
sequence=DDRRUUUULLDRURD
error=
\end{lstlisting}

\begin{table}[h]
\centering
\small
\begin{tabular}{@{}lll@{}}
\toprule
\textbf{Champ} & \textbf{Signification} & \textbf{Usage futur BDD} \\
\midrule
\texttt{level\_id} & Identifiant & Cle \texttt{UPDATE} \\
\texttt{solvable} & \texttt{yes} / \texttt{no} & Validation niveau \\
\texttt{mode} & \texttt{safe-path} ou \texttt{gems-only} & \texttt{solution\_safe} (1 si safe) \\
\texttt{move\_count} & Nombre de glissements & Statistiques \\
\texttt{sequence} & Chaine \texttt{UDLR} & \texttt{solution\_cache} \\
\texttt{error} & Message d'echec & Logs \\
\bottomrule
\end{tabular}
\caption{Champs du fichier solution}
\end{table}

\subsection{Format niveau (entree, colonne \texttt{map})}

\begin{lstlisting}[style=kllisting]
W 9
H 7
P 1 1
R 7 9
MAP
#########
#.......#
...
\end{lstlisting}

Symboles : \texttt{\#} mur, \texttt{.}/\texttt{o}/\texttt{c} collectibles, \texttt{*} portail, \texttt{\_} vide.

%------------------------------------------------------------
\section{Modes de resolution C}
%------------------------------------------------------------

\begin{table}[h]
\centering
\small
\begin{tabular}{@{}llp{6.5cm}@{}}
\toprule
\textbf{Mode} & \textbf{Fonction} & \textbf{Fantomes} \\
\midrule
safe-path & \texttt{kl\_solve\_safe()} & R, V, B pre-calcules ; jaune \textbf{ignore} \\
gems-only & \texttt{kl\_solve()} & Aucun \\
fallback & safe puis gems-only & Comme ci-dessus \\
\bottomrule
\end{tabular}
\end{table}

\textbf{Limite :} maximum \textbf{30 gemmes} (masque \texttt{uint32\_t}). Cap memoire BFS : environ 8 millions de noeuds (defini dans \texttt{solver.c}).

\textbf{Ecart avec le Worker JS :} le solveur navigateur inclut le fantome \textbf{jaune} dans l'etat BFS ; le C en mode safe-path ne le fait pas encore. Les solutions \texttt{safe-path} C peuvent donc differer du jeu sur les niveaux avec \texttt{Y}.

%------------------------------------------------------------
\section{Procedure d'utilisation}
%------------------------------------------------------------

\subsection{Compilation (une fois)}

\textbf{Windows (MinGW / MSYS2 dans le PATH) :}
\begin{lstlisting}[style=kllisting]
cd c:\MAMP\htdocs\project\solveurCrenvoyanttxt
build.bat
\end{lstlisting}

\textbf{Linux / WSL / macOS :}
\begin{lstlisting}[style=kllisting]
cd solveurCrenvoyanttxt
make
\end{lstlisting}

Produit : \texttt{solver.exe} et \texttt{solve\_db.exe} (ou sans extension sous Unix).

\subsection{Execution du pipeline complet}

\begin{enumerate}
    \item Demarrer MySQL (MAMP).
    \item Verifier \texttt{web/includes/config.php}.
    \item S'assurer que la table \texttt{niveau} contient des niveaux.
\end{enumerate}

\begin{lstlisting}[style=kllisting]
php c:\MAMP\htdocs\project\solveurCrenvoyanttxt\cli\solve_from_db.php
\end{lstlisting}

Variantes :
\begin{lstlisting}[style=kllisting]
php solveurCrenvoyanttxt/cli/solve_from_db.php --gems-only
php solveurCrenvoyanttxt/cli/solve_from_db.php --no-fallback
\end{lstlisting}

\subsection{Utilisation manuelle du binaire C}

\begin{lstlisting}[style=kllisting]
solve_db --manifest output/levels_manifest.txt --output-dir output/solutions --safe-path --fallback
\end{lstlisting}

%------------------------------------------------------------
\section{Arborescence ajoutee}
%------------------------------------------------------------

\begin{lstlisting}[style=kllisting]
project/
|-- solveurCrenvoyanttxt/
|   |-- src/
|   |   |-- batch_io.c, batch_io.h
|   |   |-- solve_db.c, solve_one.c
|   |-- cli/solve_from_db.php
|   |-- output/levels/, output/solutions/
|   |-- Makefile, build.bat, README.md
|   `-- solve_db.exe (genere)
`-- solver/                        (CLI unitaire, inchange)
    |-- solver.c, solver.h, main.c
    `-- Makefile
\end{lstlisting}

%------------------------------------------------------------
\section{Lien avec la base de donnees}
%------------------------------------------------------------

Schema pertinent (\texttt{sql/schema.sql}) :

\begin{lstlisting}[style=kllisting]
CREATE TABLE `niveau` (
  `id` INT(11) NOT NULL AUTO_INCREMENT,
  `difficulte` INT(11) NOT NULL DEFAULT 1,
  `score_max` INT(11) NOT NULL DEFAULT 0,
  `map` TEXT NOT NULL,
  `solution_cache` TEXT NULL,
  `solution_safe` TINYINT(1) NOT NULL DEFAULT 0,
  PRIMARY KEY (`id`)
);
\end{lstlisting}

\textbf{Realise :} lecture de \texttt{id} et \texttt{map} uniquement.

\textbf{Non realise (traitement ulterieur demande par l'utilisateur) :}
\begin{itemize}
    \item \texttt{UPDATE} automatique de \texttt{solution\_cache} / \texttt{solution\_safe}
    \item Modification de \texttt{get\_solution.php} ou \texttt{game.js}
    \item Suppression du Worker JavaScript
\end{itemize}

%------------------------------------------------------------
\section{Comparaison C vs JavaScript (rappel)}
%------------------------------------------------------------

Pour le \textbf{calcul BFS} pur, le C compile (\texttt{-O2}) est en general \textbf{plus rapide} que JavaScript (souvent 2x a 10x selon la taille du graphe). Une fois les solutions en cache BDD, l'affichage en jeu est \textbf{instantane} dans les deux cas.

\begin{table}[h]
\centering
\small
\begin{tabular}{@{}lll@{}}
\toprule
\textbf{Critere} & \textbf{C (batch)} & \textbf{Worker JS} \\
\midrule
Debit BFS & Eleve (natif) & Moyen (V8) \\
Fantome jaune & Non (safe C) & Oui \\
Usage en jeu & Via cache BDD & Direct ou cache \\
\bottomrule
\end{tabular}
\end{table}

%------------------------------------------------------------
\section{Limites et points d'attention}
%------------------------------------------------------------

\begin{enumerate}
    \item \textbf{Compilation requise} --- sans \texttt{solve\_db.exe}, PHP affiche une erreur avec les instructions \texttt{build.bat} / \texttt{make}.
    \item \textbf{Fantome jaune} --- ecart possible entre solution C safe-path et moteur de jeu / Worker JS.
    \item \textbf{Niveaux tres difficiles} --- risque d'epuisement des 8M noeuds BFS (pas de timeout ajoute dans le batch).
    \item \textbf{Securite} --- \texttt{solve\_from\_db.php} est un script CLI ; ne pas l'exposer en URL publique.
    \item \textbf{Environnement teste} --- \texttt{gcc} et \texttt{make} n'etaient pas dans le PATH Windows au moment du developpement ; compilation a valider localement.
\end{enumerate}

%------------------------------------------------------------
\section{Evolutions prevues (non codees)}
%------------------------------------------------------------

\begin{itemize}
    \item Script \texttt{cli/import\_solutions.php} : lire les \texttt{.txt} $\rightarrow$ \texttt{UPDATE niveau}.
    \item Porter le fantome jaune dans \texttt{kl\_solve\_safe()} pour aligner C et JS.
    \item Retirer \texttt{solver-worker.js} une fois tout importe en BDD.
    \item Tache planifiee apres ajout d'un niveau en administration.
\end{itemize}

%------------------------------------------------------------
\section{Synthese en une phrase}
%------------------------------------------------------------

\begin{quote}
Un script PHP exporte tous les \texttt{map} de la table \texttt{niveau}, le programme C \texttt{solve\_db} les resout et produit un fichier texte structure par niveau dans \texttt{output/solutions/}, pret pour un import manuel ou automatise en base --- sans modifier le site web ni le solveur JavaScript a ce stade.
\end{quote}

\vfill
\hrule
\small
\noindent
\textbf{Fichier source LaTeX :} \texttt{docs/recap\_solveur\_c\_bdd.tex} --- Projet Les fantômes d'Ombrequatre, MAMP \texttt{htdocs/project}.

\end{document}
